\documentclass[book.tex]{subfiles}

\begin{document}
\chapter{Steerable Equivariant Graph Neural Networks}

\label{chap:steerable}
\noindent\rule{11cm}{0.4pt} \\
\textbf{Prerequisites:} \autoref{chap:qft_basics}  \\
\textbf{Difficulty Level:} ** \\
\noindent\rule{11cm}{0.4pt}
\vspace{5mm}

%"""
%Including covariant information, such as position, force, velocity or spin is important in many tasks in computational physics and chemistry. We introduce Steerable E(3) Equivariant Graph Neural Networks (SEGNNs) that generalise equivariant graph networks, such that node and edge attributes are not restricted to invariant scalars, but can contain covariant information, such as vectors or tensors. This model, composed of steerable MLPs, is able to incorporate geometric and physical information in both the message and update functions. Through the definition of steerable node attributes, the MLPs provide a new class of activation functions for general use with steerable feature fields. We discuss ours and related work through the lens of equivariant non-linear convolutions, which further allows us to pin-point the successful components of SEGNNs: non-linear message aggregation improves upon classic linear (steerable) point convolutions; steerable messages improve upon recent equivariant graph networks that send invariant messages. We demonstrate the effectiveness of our method on several tasks in computational physics and chemistry and provide extensive ablation studies.
%"""

Implementing an equivariant architecture often requires formidable constructions involving spherical harmonics and Clebsch-Gordon products. The theory of steerability \cite{cohen2016steerable} offers a powerful tool to allow for simpler construction of equivariant convolutions for certain special cases. In this chapter, we use this theory to construct generable steerable equivariant graph neural networks that provide a powerful and general equivariant algorithmic framework.

%Equivariant graph neural networks in their simplest form are limited to computing invariant scalar quantities, but for more complex applications, it can be very useful to maintain more complex covariant quantities such as vectors or tensors. In this chapter, we introduce steerable $E(3)$-equivariant graph neural networks \cite{brandstetter2021geometric}, \cite{weiler2019general}, \cite{cohen2016steerable},  \cite{weiler2018learning}, \cite{weiler20183d} which maintain higher order covariant information. The architecture depends on $E(3)$-equivariant message passing layers which implement a nonlinear group convolution operation. 

%The crucial building block of this architecture is the steerable multilayer perceptrons. The architecture composes steerable MLPs along with message and update functions for the graph neural network which use geometric and physical information. The combination of non-linear message aggregation and steerable messages allows for improvement over simpler linear steerable and graph equivariant architectures.

\section{Steerable Features}

\begin{figure}
    \centering
    \includegraphics{figures/Differentiable Physics/Molecular ML/Steerable GNNs/steerable_features.png}
    \caption{\url{https://arxiv.org/abs/2110.02905}}
    \label{fig:my_label}
\end{figure}
Let 
\begin{align}\Phi: \mathcal{F} \to \mathcal{F'}\end{align}
be a convolutional network. Let \begin{align}(\mathcal{F}, \pi)
\end{align}
be a feature space with a group representation $\pi$. We say that linear space $\mathcal{F'}$ is linearly steerable with respect to $G$ if for all $g \in G$, we have that features $\Phi f$, $\Phi \pi(g) f$ are related by a linear transformation $\pi'(g)$ which does not depend on $f$. That is,
\begin{align}
\pi'(g) \Phi f &= \Phi \pi(g) f
\end{align}

This means that $\pi'(g)$ allows the "steering" of features in $\mathcal{F'}$ without referring to $\mathcal{F}$. We define a steerable vector space to be a vector space that is steerable with respect to some group. 

Using these definitions, we have $\pi'$ must also be a group representation
\begin{align}
    \pi'(gh)\Phi f &=  \Phi \pi(gh) \\
    &= \Phi \pi(g) \pi(h) f \\
    &= \pi'(g) \Phi \pi(h) f \\
    &= \pi'(g)\pi'(h) \Phi f
\end{align}
The intuition here is that steerable networks allow for linear transformations to implement the effect of the group action. $SO(3)$ for example naturally allow for steerable transformations since any rotation matrix $R$ transforms a vector $v$ by multiplication as $Rv$. An arbitrary group action however may not be steerable. 

What is the difference between equivariance and steerability in general? 
Steerable vector spaces are constructed by spherical harmonic spaces, but not all equivariant methods use steerability as a guiding principle.

\section{Equivariant Messages for Graph Networks}

Geometrical and physical quantities improve equivariant mesage passing by allowing message passing networks to gather vectorial and tensorial messages. Steerable equivariant graph neural networks provide a family of techniques for passing equivariant messages in graph neural networks to allow for construction of a general family of methods.

How do we build steerable $E(3)$ equivariant layers? Note that for $g \in SO(3)$ and vector $\tilde{h}$ in an irreducible representation, we have that the group action is given by
\begin{align}
D^\ell(g)\tilde{h},\ g \in SO(3)
\end{align}
Here $D^\ell$ is a Wigner $D$-matrix. As a special case, a Euclidean vector in $\mathbb{R}^3$ is consequently steerable for rotations $g = R \in SO(3)$ via $D^1(g) = R$.

Of course, we don't only want to work with 3D vectors. By choosing larger values of $\ell$, we can have larger equivariant message passing spaces. In designing such a generalized steerability, we would like to meet two design goals.

\begin{itemize}
\item Goal 1: Generalization of notion of rotations to arbitrary large vector spaces
\item Goal 2: $SO(3)$-equivariance
\end{itemize}

\section{Constructing Suitable Equivariant Message Passing Spaces}

We use an orthonormal basis for functions on the sphere where we represent any vector in a basis spanned by spherical harmonics.
\begin{align}
    h \in V_L = V_0 \oplus V_1 \oplus V_2 \oplus \dotsc
\end{align}

The $V_l$ are steerable subspaces since $SO(3)$ has a group representation via Wigner $D$-matrices $D^{(l)}(g)$ which act on each subspace $V_l$ separately. $D^{(l)}(g)$ are block diagonal representations with each block acting on a separate vector space. This  property achieves Goal 1 of processing high dimensional vectors.

\begin{figure}
    \centering
    \includegraphics{figures/Differentiable Physics/Molecular ML/Steerable GNNs/steerable_vector.png}
    \caption{\url{https://arxiv.org/abs/2110.02905}}
    \label{fig:my_label}
\end{figure}


To construct more general steerable spaces, we use a spherical harmonics embedding. Recall the spherical harmonics are functions defined on the sphere
\begin{align}
    Y^{(l)}_m: S^2 \to \mathbb{R}
\end{align}

We can evaluate a spherical harmonic for an arbitrary vector $x$ in $\mathbb{R}^3$ by normalizing the vector
\begin{align}
Y^{(l)}_m\left (\frac{x}{\|x\|}\right)
\end{align}
By ranging over all the possible choices of $m$, we can construct an embedding vector $a^{(l)}(x)$
\begin{align}
    a^{(l)}(x) := (Y^{(l)}_{-l}\left (\frac{x}{\|x\|}\right),\dotsc,Y^{(l)}_l\left (\frac{x}{\|x\|}\right))^T
\end{align}
This vector is called the spherical harmonics embedding of the vector $x \in \mathbb{R}^3$. We note that the spherical harmonics embedding is $SO(3)$ equivariant due to the equivariance of the spherical harmonics. That is, we have the property that 
\begin{align}
    \forall R' \in SO(3), \forall n \in S^2, a^{(l)}(R'n) = D^{(l)}(R')a^{(l)}(n)
\end{align}
This equivariance allows us to achieve Goal 2 formulated above.




\section{Steerable Multilayer Perceptron}
Steerable multilayer perceptrons act like regular multilayer perceptrons by interleaving nonlinearities with linear equivariant transformations. Consequently, we want equivariant linear maps (also known as intertwiners) which transform between steerable vector spaces at layer $i-1$ and layer $i$ by a linear transformation
\begin{align}
    h^i = W^i_{\tilde{a}} h^{i-1}
\end{align}
Unlike regular MLPs, the weight matrices $W^i$ must be constructed in a structured fashion using Clebsch-Gordon products as in tensor field networks or other equivariant architectures we have seen thus far.

Recall that the Clebsch-Gordon tensor product gives the linear map
\begin{align}
    (h^{l_1} \otimes^w_{cg} h^{(l_2)})^{(l)}_m = w \sum_{m_1=-l_1}^{l_1} \sum_{m_2 = -l_2}^{l_2} C^{(l,m)}_{(l_1,m_1),(l_2, m_2)} h^{(l_1)}_{m_1}h^{(l_2)}_{m_2}
\end{align}
We treat the Clebsch-Gordon product with a fixed vector $\tilde{a}$ in one of its inputs as a steerable linear layer conditioned on $\tilde{a}$
\begin{align}
W_{\tilde{a}}\tilde{h}^{i-1} = \tilde{h}^{i-1} \otimes_{cg}^{w} \tilde{a}
\end{align}

If the nonlinearity is chosen to be steerable, the resulting steerable MLP will satisfy the property

\begin{align}
    \tilde{MLP}(D(g)\tilde{h}_0) = D'(g) \tilde{MLP}(\tilde{h}_0)
\end{align}

There are few choices of steerable nonlinearity such as gated nonlinearities.

\section{Steerable Equivariant Graph Neural Networks}

As in our earlier discussion of $SE(3)$ transformers, point convolutions are given by the equation
\begin{align}
    f' = \sum_{j \in \mathcal{N}(i)} W(x_j - x_j) f_j
\end{align}
Messages are given by
\begin{align}
    m_{ij} = W(x_j x_i)f_j
\end{align}
The message update rule is
\begin{align}
    f'_i = \sum_{j\in\mathcal{N}(i)} m_{ij}
\end{align}
The message passing rule for steerable equivariant graph neural networks is given by
\begin{align}
    m_{ij} &= \phi_m(f_i,f_j,\|x_j-x_i\|^2), a_{ij}) \\
    f'_i &= \phi_f(f_i, \sum_{j \in \mathcal{N}(i)} m_{ij} a_i) \\
    x'_i &= x_i + C \sum_{j\neq i} (x_i - x_j) \phi_x(m_{ij})
\end{align}
This is a nonlinear convolution with an anisotropic kernel. Here the $a_{ij}, a_i$ are steerable features for edges and nodes respectively. The $\phi_m$ and $\phi_f$ are steerable MLPs. The edge attributes $\tilde{a}_{ij}$ could be for example relative position, relative force, and the node attributes $\tilde{a}_i$ could be velocity, force, spin, etc. This structure allows for directly leveraging local geometric or physical cues in the node update, which leads to a new class of steerable activation functions.
\begin{align}
    m_{ij} &= MLP(f_i, f_j, \|x_j-x_i\|^2, a_i) = \sigma(W^{(k)}(\dotsc(\sigma(W^{(1)}h_{i-1})))) \\
    h_{i-1} &= f_i || f_j || \|x_j - x_i\|^2 || a_i
\end{align}

Steerable kernel methods expand the linear transformations in a steerable basis such as 3D spherical harmonics
\begin{align}
    W(x_j-x_i) = \sum_l\sum_{m=-l}^l W^{(l)_m}(\|x_j - x_i\|) Y^{(l)}_m(x_j - x_i)
\end{align}
Equivariant steerable methods can be written in convolution form
\begin{align}
    \tilde{f}'_j = \sum_{j \in \mathcal{N}(i)} W_{\tilde{a}_{ij}}(\|x_j - x_i\|)\tilde{f}_j
    \tilde{f}'_j = \sum_{j \in \mathcal{N}(i)} W_{\tilde{a}_{ij}}(\tilde{f}_i, \tilde{f}_j, \|x_j - x_i\|)\tilde{f}_j
\end{align}
Note that this same transformation is used in tensor field networks, Cormorant, L1Net \cite{miller2020relevance} and the $SE(3)$ transformer. We can then look at equivariant message passing as a nonlinear convolution. 


%A steerable function is said to have steerable output for all input if the transformed input and non-transformed output are only differentiable by something that doesn't depend on the input.

%We can then think about how to bring steerability into other layers.

%\begin{itemize}
%\item messages are non-linear transformations
%\end{itemize}

%SEGNNs
%\begin{align}
%    m_{ij} = \tilde{MLP}(\tilde{f}_i, \tilde{f}_j, \|x_j - x_i\|^2, \tilde{a}_{ij})
%    \tilde{f}'_i = \tilde{MLP}(\tilde{f}_i, \sum_{j \in \mathcal{N}(i)} \tilde{m}_{ij}, \tilde{a}_j
%\end{align}



%\section{Physical Applications}

%One potential use of a SEGNN is a N-body experiment. Here assume you have 5 particles with initial position, veloccity and charge. We want to predict position after 1000 timesteps. Input $h \in V_0 \oplus V_1 \oplus V_1$ (relative position to center of mass, velocity, norm of velocity) so 2 vectors + 1 scalar = 7 dimensional.

%Output is $\tilde{o} \in V_1$ is a vector quantity.

%Tensor Field Networks here use linear messages.


%Nonlinear outperforms linear convolutions. So far CNNs mostly use linear convolutions.
\end{document}